\section*{Capítulo \ref{ch:b3:quantum-angular-momentum} --- Momento angular cuántico}

\begin{solution}{exo:b3:quantum-angular-momentum:1}
(a) $[\hat y\hat p_z - \hat z\hat p_y,\ \hat z\hat p_x - \hat x\hat
p_z]$: solo sobreviven los términos que comparten una pareja $\hat z, \hat p_z$,
lo que da $\iu\hbar(\hat x\hat p_y - \hat y\hat p_x)$. (b) $[\hat
L^2, \hat L_z] = \sum_i[\hat L_i^2, \hat L_z] = \hat L_x[\hat L_x,
\hat L_z] + [\hat L_x, \hat L_z]\hat L_x + (x \to y)$: los cuatro
términos se cancelan por parejas. (c) $[\hat L_z, \hat x] = \iu\hbar\hat y$:
$\hat L_z$ genera las rotaciones en torno a $z$ --- tanto de operadores como
de estados. (d) Que dos componentes estén simultáneamente bien definidas
obliga (por el \omterm{def:b3:quantum-formalism:commutator}{conmutador}) a que la tercera también lo esté con valor nulo,
y con ella las tres: solo $l = 0$ lo consigue.
\end{solution}

\begin{solution}{exo:b3:quantum-angular-momentum:2}
(a) $\hat L_z = \hbar\operatorname{diag}(1, 0, -1)$. (b) $\hat L_+ =
\hbar\sqrt2\,\begin{pmatrix}0&1&0\\0&0&1\\0&0&0
\end{pmatrix}$, $\hat L_-$ su traspuesta. (c) $\hat L_x =
\dfrac{\hbar}{\sqrt2}\begin{pmatrix}0&1&0\\1&0&1\\0&1&0
\end{pmatrix}$: polinomio característico $\lambda(\lambda^2 -
\hbar^2)$. (d) El vector propio de $L_x = +\hbar$ es $\tfrac12(1,
\sqrt2, 1)$: probabilidades $\tfrac14, \tfrac12, \tfrac14$ para $m =
+1, 0, -1$.
\end{solution}

\begin{solution}{exo:b3:quantum-angular-momentum:3}
(a) $I = \mu r_0^2 = \qty{1.46e-46}{kg.m^2}$; $B = \hbar^2/2I =
\qty{3.8e-23}{J} = \qty{0.24}{meV}$. (b) $2B/h = \qty{115}{GHz}$,
$\lambda = \qty{2.6}{mm}$. (c) \qty{231}{GHz} y \qty{346}{GHz}.
(d) El espaciado da $B$ y, con él, $I = \mu r_0^2$: la longitud de enlace de
una molécula situada a años luz, con aproximadamente la mitad de la
precisión relativa del espaciado.
\end{solution}

\begin{solution}{exo:b3:quantum-angular-momentum:4}
(a) $m\hbar$ con $m = -2 \dots 2$. (b) $\cos\theta = 2/\sqrt6$:
$\theta = 35.3^\circ$. (c) Un alineamiento perfecto haría que $L_x = L_y
= 0$ estuvieran bien definidos a la vez que $L_z$ --- algo que prohíben los
\omterm{def:b3:quantum-formalism:commutator}{conmutadores} salvo en el caso trivial $l = 0$. (d) $l/\sqrt{l(l+1)} \to 1$:
para $l$ grande el cono se cierra sobre el eje y vuelve la flecha clásica.
\end{solution}

\begin{solution}{exo:b3:quantum-angular-momentum:5}
(a) Sin dependencia en $\varphi$, el operador da
$-\hbar^2(\sin\theta)^{-1}\partial_\theta(\sin\theta\,(-\sin\theta))
= 2\hbar^2\cos\theta$. (b) El mismo cálculo con el factor
$\eu^{\pm\iu\varphi}$: \omterm{def:b3:quantum-formalism:observable}{valor propio} $2\hbar^2$; $\hat L_z$
da $\pm\hbar$. (c) $\cos\theta$ y $\sin\theta\,\eu^{\pm\iu
\varphi}$ cambian de signo bajo $\vect r \to -\vect r$ ($\theta \to \pi
- \theta$, $\varphi \to \varphi + \pi$): paridad $-1 = (-1)^1$. (d)
Son vectores propios del \omterm{def:b3:quantum-formalism:observable}{operador hermítico} $\hat L_z$ con \omterm{def:b3:quantum-formalism:observable}{valores propios}
distintos.
\end{solution}

\begin{solution}{exo:b3:quantum-angular-momentum:6}
(a) $2J + 1$ estados comparten el nivel (degeneración); el factor de
Boltzmann grava su energía. (b) Maximícese el producto: $\dd/\dd J = 0$
da $2 = (2J+1)^2B/k_{\text{B}}T$, el $J_{\max}$ enunciado. (c) A
\qty{10}{K}: $J_{\max} \approx 0.8$, de modo que domina $J = 1$ y brillan
las líneas de \qty{115}{GHz} y \qty{230}{GHz}; a \qty{300}{K}:
$J_{\max} \approx 7$. (d) $T \approx 2B(J_{\max} + \tfrac12)^2/
k_{\text{B}} \approx \qty{310}{K}$ --- un termómetro rotacional.
\end{solution}

\begin{solution}{exo:b3:quantum-angular-momentum:7}
(a) El nivel superior se desdobla en cinco y el inferior en tres, todos
con el mismo espaciado $\mu_{\text{B}}B$. (b) Con espaciados iguales,
$h\nu = h\nu_0 + (\Delta m)\mu_{\text{B}}B$ y $\Delta m \in \{0,
\pm1\}$: solo tres posiciones. (c) $\mu_{\text{B}}/h = \qty{14}{GHz/T}$, de modo que a $B = \qty{2}{T}$
el desdoblamiento vale \qty{28}{GHz}, es decir, $\Delta\lambda =
\lambda^2\Delta\nu/c \approx \qty{23}{pm}$ --- resoluble desde las redes de
difracción de Zeeman de 1896.
(d) El espín propio del electrón, con su factor anómalo $g \approx
2$: los patrones ``anómalos'' fueron las huellas dactilares del espín antes
de que el espín tuviera nombre (\cref{ch:b3:spin-two-level}).
\end{solution}

\begin{solution}{exo:b3:quantum-angular-momentum:8}
En las unidades enunciadas, $U_{\text{ef}} = -2/x + l(l+1)/x^2$. (a) Para
$l = 1$: mínimo en $x = 2$ ($r = 2a_0$), de profundidad $-E_{\text{I}}/2$.
(b) $U_{\text{ef}} > 0$ para $x < 1$: dentro de un radio de Bohr gana la
muralla. (c) $\psi \sim r^l$ cerca del origen: solo los estados $l = 0$
tienen densidad no nula en $r = 0$. (d) Capturar un electrón exige densidad
electrónica \emph{en} el núcleo: los electrones $s$, los únicos no vedados
por la muralla centrífuga, son los que se tragan.
\end{solution}

\begin{solution}{exo:b3:quantum-angular-momentum:9}
(a) $\hat L_x = (\hat L_+ + \hat L_-)/2$ cambia $m$: los elementos
diagonales se anulan. (b) Por simetría, los dos cuadrados transversales son
iguales, y su suma vale $\langle\hat L^2 - \hat L_z^2\rangle$. (c)
En $m = l$: $\Delta L_x\Delta L_y = l\hbar^2/2 = \tfrac\hbar2
\,l\hbar$: igualdad --- el estado estirado está tan alineado como el
álgebra permite. (d) El cono: $L_z$ bien definido, medias transversales
nulas y dispersiones transversales iguales --- un vector de longitud y
proyección definidas, difuminado democráticamente en acimut.
\end{solution}

\begin{solution}{exo:b3:quantum-angular-momentum:10}
(a) Contabilidad energética con $E_J = BJ(J+1)$: $J \to J + 1$ suma
$2B(J+1)$; $J \to J - 1$ resta $2BJ$. (b) $\Delta J = 0$ está
prohibido, y las dos ramas arrancan en $\pm 2B$: nada cae en
la frecuencia vibracional pura. (c) Dos peines de espaciado $2B$
que flanquean un hueco, con intensidades que crecen hasta $J_{\max}$ y luego
caen. (d) La anchura de la banda es la extensión poblada de los peines, que
crece como $\sqrt{T}$: esas alas pobladas térmicamente son precisamente
donde una banda de invernadero saturada sigue absorbiendo.
\end{solution}

\begin{solution}{exo:b3:quantum-angular-momentum:11}
(a) $\|\hat J_\pm\ket{j,m}\|^2 = \bra{j,m}\hat J_\mp\hat J_\pm
\ket{j,m} = \hbar^2[j(j+1) - m(m\pm1)]$. (b) Subir más allá de $m = j$
daría $j(j+1) - j(j+1) = 0$ y, un paso más allá, normas negativas:
la cadena debe contener el estado superior aniquilado. (c) La
cima y el fondo difieren en un número entero de pasos unidad: $2j \in
\mathbb N$: $j = 0, \tfrac12, 1, \tfrac32, \dots$ (d) El álgebra
no usa funciones de onda en ningún momento; solo la exigencia de que
$\eu^{\iu m\varphi}$ sea univaluada sobre la circunferencia obliga a que $m$
sea entero --- una condición que ata al movimiento orbital y deja libres a
los momentos angulares intrínsecos (de espín) para ser semienteros.
\end{solution}

\begin{solution}{exo:b3:quantum-angular-momentum:12}
(a) $n\hbar$ por unidad de volumen. (b) $\omega = n\hbar V/(\tfrac12 Ma^2)
= 2n\hbar/\rho a^2 = 2 \times \num{8.5e28} \times \num{1.055e-34}/
(7900 \times \num{e-4}) \approx \qty{2.3e-5}{rad/s}$ --- invisible
de una sola vez. (c) Invertir al compás de la resonancia del hilo
acumula $\sim Q$ impulsos: $\omega \sim \qty{2e-3}{rad/s}$,
oscilaciones de milirradián con un periodo de segundos --- visibles con un
espejo y un haz de luz, como vieron Einstein y de Haas. (d) La razón
giromagnética medida fue $e/m_{\text{e}}$, el doble de la orbital
$e/2m_{\text{e}}$: el magnetismo del hierro no son corrientes orbitales,
sino el \emph{espín} del electrón, cuyo $g \approx 2$ explican los capítulos
siguientes.
\end{solution}

\begin{solution}{pb:b3:quantum-angular-momentum:1}
\textbf{1.} $E_J = BJ(J+1)$, con degeneración $2J + 1$.
\textbf{2.} $I = \mu r_0^2 = \qty{1.46e-46}{kg.m^2}$: $B =
\hbar^2/2I = \qty{3.8e-23}{J}$, es decir, $B/h = \qty{57.7}{GHz}$ ---
la constante medida.
\textbf{3.} $115$, $231$ y \qty{346}{GHz}.
\textbf{4.} $E_{J+1} - E_J = 2B(J+1)$ crece linealmente: líneas
consecutivas difieren en la constante $2B$. El espaciado da $I$ y, con él,
la longitud de enlace --- química estructural por radio.
\textbf{5.} $\lambda = c/\nu = \qty{2.6}{mm}$. El vapor de agua
atmosférico absorbe con avidez las ondas milimétricas; de ahí Atacama, Mauna
Kea o el Polo Sur --- altos, fríos y secos.
\textbf{6.} El fotón es una partícula de espín $1$ que lleva $\hbar$:
la conservación del momento angular total obliga al $J$ de la molécula a
cambiar en una unidad.
\textbf{7.} Su distribución de carga es simétrica para cualquier ángulo de
rotación: sin dipolo oscilante no hay línea dipolar --- camuflaje
perfecto.
\textbf{8.} $B_{\text{H}_2} = \hbar^2/2\mu r_0^2 \approx
\qty{7.6}{meV}$; las débiles transiciones cuadrupolares de la molécula
simétrica exigen $\Delta J = 2$, con un coste de $6B \approx
\qty{46}{meV}$ --- cincuenta veces $k_{\text{B}}T$ a \qty{10}{K}: el
hidrógeno frío no puede radiar en absoluto.
\textbf{9.} $2B = \qty{0.48}{meV} < k_{\text{B}}T = \qty{0.86}{meV}$:
las colisiones mantienen poblados los primeros peldaños --- la escalera
funciona justo donde viven las nubes.
\textbf{10.} $n_1/n_0 = 3\,\eu^{-0.478/0.862} = 1.7$: el nivel emisor está
bien surtido.
\textbf{11.} $J_{\max} \approx \sqrt{0.86/0.48} - 0.5 \approx 0.8$:
la población tiene su máximo en $J = 1$.
\textbf{12.} El CO entrega los fotones; convertirlos en masa total de gas
descansa en una abundancia supuesta de CO respecto del H$_2$ --- mensajero
luminoso, rescate calibrado.
\textbf{13.} $\Delta\nu = \qty{115}{MHz}$: $v = c\,\Delta\nu/\nu
\approx \qty{300}{km/s}$, alejándose (la frecuencia ha bajado) ---
velocidades orbitales galácticas.
\textbf{14.} Los $\pm\qty{2}{km/s}$ abarcan \qty{1.5}{MHz}; el movimiento
térmico a \qty{10}{K} da solo $\sim\qty{40}{kHz}$: la anchura es
turbulencia, no temperatura --- las anchuras de línea cartografían la
meteorología interna de una nube.
\textbf{15.} En torno a una masa central, $v \propto 1/\sqrt r$ debería
caer (Kepler); la planitud medida significa que la masa encerrada sigue
creciendo con el radio --- materia invisible que envuelve al disco
luminoso: materia oscura.
\textbf{16.} La razón de poblaciones, es decir, la temperatura de
excitación del gas.
\textbf{17.} A \qty{57.6}{GHz}: la expansión cósmica ha duplicado toda
longitud de onda en vuelo --- el corrimiento al rojo cosmológico al que
vuelve el último capítulo de este libro.
\textbf{18.} Flujo de fotones $\to$ luminosidad de CO (hace falta la
distancia) $\to$ recuento de CO (modelo de excitación) $\to$ masa de H$_2$
(el ``factor X'' de CO a H$_2$: el eslabón más débil, calibrado localmente
y exportado con prudencia).
\textbf{19.} Un núcleo a \qty{10}{K} no emite nada que pueda ver un
telescopio óptico; sus líneas rotacionales son a la vez su único
termómetro, su velocímetro y su balanza.
\textbf{20.} Un termómetro es sensible donde el espaciado de niveles es $\sim
k_{\text{B}}T$: un espaciado de \qty{5.5}{K} responde con fuerza
entre $5$ y $50$ K, mientras que un nivel óptico de \qty{2}{eV} no lo
poblaría absolutamente nada.
\textbf{21.} La línea se satura (se vuelve ópticamente gruesa) y, peor aún,
el CO se congela sobre los granos de polvo en el gas más denso y frío: los
astrónomos pasan entonces a isotopólogos más raros ($^{13}$CO, C$^{18}$O)
y a otras moléculas.
\textbf{22.} Cada especie necesita su propia densidad y temperatura para
brillar: juntas hacen una tomografía de la nube --- las envolturas
exteriores en CO, los núcleos densos en NH$_3$ y HCN.
\textbf{23.} Una línea se ve por contraste frente al fondo cósmico; cuando
la temperatura del gas se acerca a \qty{2.7}{K}, su excitación se equilibra
con el fondo y el contraste se desvanece: ninguna nube más fría puede leerse
así.
\textbf{24.} $\lambda/20 = \qty{130}{\micro m}$ frente a
$\qty{25}{nm}$ para el trabajo óptico: cinco mil veces más
indulgente --- por eso las antenas milimétricas de doce metros están al aire
libre en el desierto mientras los espejos ópticos viven en cúpulas.
\textbf{25.} Una escalera $BJ(J+1)$ con $B/h = \qty{57.6}{GHz}$, viva a
diez kelvin y leída en \qty{2.6}{mm}, mide la masa, la
temperatura, la turbulencia y la rotación del universo frío --- y sus
curvas de rotación planas son una exposición permanente sobre la materia
oscura.
\end{solution}
